\begin{appendix}
\section{Bivariate renewal theory, important estimates}\label{app:bivrenewal}

We present here some results on the bivariate renewal process $\tau$ defined in Section~\ref{sec:model}, and in particular Proposition~\ref{prop:transience} which gives some conditions on the transience/recurrence of the intersection renewal $\sigma=\tau\cap\tau'$. Recall the notations of Section~\ref{sec:defan} for the recentering sequence $(b_n)_{n\geq 0}$ and for the scaling sequence $(a_n)_{n\geq 0}$.


\subsection{Local large deviations and a useful Lemma}

We first present some local large deviation estimate, which is used in the proof of Lemma~\ref{lem:homogen}.

\begin{theo}[{\cite[Theorem~2.4]{cf:B}}]\label{thm:locallargedev}
Assume that $\mu<+\infty$. We have that there exists a constant $C>0$ such that uniformly for $(n,m)$ such that $n - \mu k \geq a_k \wedge (C\sqrt{k\log k}),$
\[
\bP(\tau_k =(n,m)) \leq C\,k\,K(|n-\mu k| + |m-\mu k|)\,.
\]
\end{theo}

We give another useful lemma, that controls the number of renewals in $(0,N]\times(0,M]$, in the case $\ga\in(0,1)$.

\begin{lemm}\label{th:limp}
Assume $\ga\in(0,1)$. Given $\delta>0$ there exists $\gep>0$ such that for $N$ sufficiently large and $M \sim \gamma N$ we have
\begin{equation}
\bP \big(\vert \tau \cap (0,N] \times (0,M] \vert \ge \gep N^{\ga}/L(N) \big) \ge 1- \delta\,.
\end{equation}
\end{lemm}

\begin{proof}
Set $n = n(\gep,N) = \gep N^{\ga}/L(N)$ and $B_{N,M} := (0,N]\times (0,M]$. We want to~prove
\begin{equation}
\bP \left(\tau_n \notin B_{N,M} \right) \le \delta\,.
\end{equation}
Let us define $\tilde \tau_n = (\tilde \tau_n^{(1)},\tilde \tau_n^{(2)})$ with
\begin{equation}
\tilde{\tau}^{(1)}_n := \sum_{i=1}^n (\tau_i^{(1)} - \tau_{i-1}^{(1)}) \unbf_{ \lbrace \tau_i^{(1)} - \tau_{i-1}^{(1)} \le N \rbrace}\,,\quad
\tilde{\tau}^{(2)}_n := \sum_{i=1}^n (\tau_i^{(2)} - \tau_{i-1}^{(2)}) \unbf_{ \lbrace \tau_i^{(2)} - \tau_{i-1}^{(2)} \le M \rbrace}\,.
\end{equation}

Then we have
\begin{equation}\label{eq:tautilde}
\bP \left(\tau_n \notin B_{N,M} \right) \le \bP \left(\tilde{\tau_n} \notin B_{N,M} \right) + \bP \left(\exists i \le n ; (\tau_i - \tau_{i-1}) \notin B_{N,M} \right).
\end{equation}
Note that the marginals $\tau^{(1)}_n$ and $\tau^{(2)}_n$ have the same distribution: as $N \to \infty$ we~have
\begin{equation}
\bP(\tau^{(1)}_1= N) = \bP(\tau^{(2)}_1= N) \sim \frac{1}{(1+\ga)} L(N) N^{-(1+ \ga)}\,.
\end{equation}

Therefore, we can bound the second term in~\eqref{eq:tautilde} by
\begin{equation}
n \left(\bP (\tau^{(1)}_1 \!>\! N)\!+\! \bP (\tau^{(2)}_1 \!>\! M) \right) \le n \left(C_{\ga} L(N) N^{-\ga} \!+\!C_{\ga} L(M) M^{-\ga}\right) \le \gep C_{\ga,\gamma}\,,
\end{equation}
which is smaller than $\gd/2$ if $\gep \leq \gd/(2 C_{\ga,\gamma})$.

The first term in~\eqref{eq:tautilde} is bounded by $ \bP(\tilde{\tau}^{(1)}_n > N) + \bP(\tilde{\tau}^{(2)}_n > M)$. Observe that for every choice of $\lambda_1 >0$ we have
\begin{equation}
\bP (\tilde{\tau}^{(1)}_n > N) \le e^{ -\lambda_1 N} \bE \big[e^{\lambda_1 \tilde{\tau}^{(1)}_n}\big]
\le e^{n \log \bE[e^{\lambda_1 \tilde{\tau}^{(1)}_1}] - \lambda_1 N }\,.\label{eq:expobound}
\end{equation}
Using the fact that $\tilde \tau_1^{(1)} = \tau_1^{(1)} \unbf_{\{\tau_1^{(1)}\leq N \}}\leq N$, we get that for any $s \ge 1$
\begin{equation}
\bE [(\tilde{\tau}^{(1)}_1)^s] \le N^{s-1}
\bE \big[\tau^{(1)}_1 \unbf_{\{\tau^{(1)}_1 \leq N\}}\big] \le c_{\ga} L(N) N^{- \ga } N^s\,,
\end{equation}
where we used that $\bP(\tau_1^{(1)} =N) \leq \cst. L(N) N^{-(1+\ga)}$ with $\ga\in(0,1)$ to estimate the second expectation. In the end, expanding the exponential and using the above bound, we get
\begin{equation}
\log \bE \left[e^{\lambda_1 \tilde{\tau}^{(1)}_1} \right] \le \log \left(1 + c_{\ga} L(N) N^{- \ga } (e^{\lambda_1 N} -1) \right) \le c_{\ga} L(N) N^{- \ga } e^{\lambda_1 N}\,.
\end{equation}

We pick $k_0$ such that $\gd^{k_0-1}\leq e^{-1}/4$, and choose $\lambda_1 = k_0 N^{-1} \log(1/ \delta)$, then with the definition of $n=\gep L(N)^{-1} N^{\ga}$, we get $n \log \bE [e^{\lambda_1 \tilde{\tau}^{(1)}_1}] \le c_{\ga} \gep \delta^{-k_0} $, and choosing $\gep\leq c_{\ga}^{-1} \gd^{k_0}$ we get from~\eqref{eq:expobound}
\begin{equation}
\bP(\tilde{\tau}^{(1)}_n > N) \le e \gd^{k_0} \leq \gd/4\,.
\end{equation}
%\pagebreak

Using the same reasoning and choosing $\lambda_2 = k_0 M^{-1} \log (1/\delta)$, we have that if $\gep\leq c_{\ga,\gamma}^{-1} \gd^{k_0}$ (for some constant $c_{\ga,\gamma}$),
\begin{equation}
\bP(\tilde{\tau}^{(2)}_n > M) \le e \delta^{k_0} \leq \gd/4\,.
\end{equation}

The proof is therefore complete by taking $\gep = \min \lbrace \delta / (2 C_{\ga,\gamma}), c_{\ga}^{-1} \delta^{k_0}, c_{\ga,\gamma}^{-1} \delta^{k_0} \rbrace$.
\end{proof}

\subsection{Renewal theorems, and the intersection of two independent copies}

The goal of this section is to estimate the mean overlap of two copies $\tau$ and $\tau'$ in the region $(0,N]\times (0,M]$. We leave aside the case $\ga=1$ which is more technical (in particular if $\mu=+\infty$): we refer to Remark~\ref{rem:alpha=1} for more comments on this case. We~define
\begin{equation}\label{eq:UNM}
U_{N,M}:= \bE\left[\left| \sigma\cap \left([0,N]\times [0,M]\right) \right|\right] =\sum_{n=0}^{N} \sum_{m=0}^{M} \bP((n,m)\in\tau)^2\,,
\end{equation}
and for any $\gl>0$
\begin{equation}
\hat U (\gl) := \sum_{n,m=0}^{+\infty} e^{-\gl (n+m)} \bP((n,m)\in\tau)^2\,.
\end{equation}

\begin{prop}\label{prop:transience}
If $\ga<1$, then $\sup_{N,M\in\bbN} \ U_{N,M} <+\infty.$ If $\ga>1$, then set $\rho:=1-\min(\ga,2)^{-1} \in [0,1/2]$. We have
\begin{equation}\label{suman}
\sum_{n=1}^N \frac{1}{a_n} \sim \varphi(N) N^{\rho} \to +\infty \qquad \text{as}\;N\to\infty,
\end{equation}
for some slowly varying function $ \varphi ({\,\cdot\,})$. Moreover,
\begin{equation}\label{eq:asympUn}
\begin{aligned}
U_{N,N}&\sim 2 c_{\ga} \varphi(N) N^{\rho} &&\quad \text{as } N\to\infty, \\
\hat U(\gl) & \sim \frac{2^{1-\rho} c_{\ga}}{\Gamma(1+\rho)} \varphi(1/\gl) \gl^{-\rho} &&\quad \text{as } \gl\downarrow 0,
\end{aligned}
\end{equation}
with $c_{\ga} = \int_{0}^{\infty} c_{\ga}(t)^2 \dd t$, $c_{\ga}(t)$ being the constant appearing in Theorem~\ref{th:reD}.
As a consequence, $\sigma=\tau\cap\tau'$ is terminating if $\ga<1$, and persistent if $\ga> 1$.
\end{prop}

This proposition is based on renewal theorems (see Theorems~\ref{th:re01}--\ref{th:reD} below), that can be found in~\cite{cf:B} (in a more general setting), giving sharp asymptotics along the favorite direction, and general upper bounds away from it. The case $\ga=1$ can also be found in~\cite{cf:B} but we do not include it here, see Remark~\ref{rem:alpha=1}.


\begin{theo}\label{th:re01}
If $\ga \in (0,1)$, then for $n\to+\infty $ and $r$ such that $r/n \to t \in \bbR_+$, we have
\begin{equation}\label{eq:D01}
\bP \left((n,n+r) \in \tau \right) \stackrel{n \to \infty}\sim C_{\ga}(t) L(n)^{-1} n^{-(2-\ga)}\,,
\end{equation}
with $C_{\ga}(t) := \ga \int_{0}^{+ \infty} x^{1- \ga} g_\ga (x, (1+t) x)\,\dd x$. Moreover, for any $\gd>0$ there is a constant $C_{\gd}>0$ such that for any $r\geq n$,
\begin{equation}
\bP\big((n,n+r)\in\tau \big) \leq C_{\gd} L(n)^{-1} n^{-(2-\ga)} \times \Big(\frac{r}{n} \Big)^{-(1+\ga) +\gd}\,.
\end{equation}
\end{theo}



\begin{theo}\label{th:reD}
If $\ga> 1$, for $n \to \infty$ and $r$ such that $r/a_n \to t \in \bbR_+$, we have that
\begin{equation}
\bP \left((n,n+r) \in \tau \right) \stackrel{n \to \infty}\sim c_{\ga}(t)\,\frac{1}{a_n}\,,
\end{equation}
where $c_{\ga}(t) = \mu_\ga \int_{- \infty}^{+ \infty} g_\ga (x,x+\mu_\ga t) \dd x$ with $\mu_\ga := \mu^{1/\min(\ga,2)}$ and $g_{\ga}$ the density of the limiting distribution of $(\tau_n -\mu n)/a_n$ ($\alpha$-stable if $\ga\in(1,2)$ or normal if $\ga\geq 2$). Moreover, for any $\gd>0$ there exists a constant $C_{\gd}>0$ such that for any $r\geq a_n$,
\begin{equation}\label{eq:offdiag}
\bP\big((n, n+r) \in \tau \big) \leq \frac{C}{a_n} \Big(\frac{r}{a_n} \Big)^{-(1+\ga) +\gd}\,.
\end{equation}
\end{theo}

Theorems~\ref{th:re01} and~\ref{th:reD} are extracted from~\cite[Theorems~3.1, 4.1 and Theorems~3.3, 4.2 respectively]{cf:B}: we refer to~\cite[Equations~(3.4), (4.2), (3.7), (4.5)]{cf:B} respectively, for a statement in the symmetric setting we are considering here.
%~question des références en dure ici ? Problème avec respectively. Ref après respectively ne sont pas des théo dans cet article mais des propositions.donc références de la biblio Ber18b citée auparavant ? Je ne sais pas...
%~ Les ref sont celles de cf:B. 
\\[-0.5em]

\begin{proofc}[Proof of Proposition~\ref{prop:transience}]

\begin{step}{Case $\ga<1$}
Notice that by symmetry, for $M\geq N$,
\[
U_{N,M} \leq U_{M,M} \leq 2 \sum_{n=1}^M \sum_{r=0}^{M-n} \bP((n,n+r) \in\tau)^2\,.
\]
We therefore need to control the last sum. Let us denote
\[
W_n:=\sum_{r\geq 0} \bP((n,n+r) \in\tau)^2\,.
\]
Using Theorem~\ref{th:re01} (and properties of slowly varying functions), we get that there is a constant $c$ such that for all $n\geq 1$
\begin{align*}
W_n \leq c\sum_{r=1}^n L(n)^{-2} n^{-2(2-\ga)} + c\sum_{r\geq n} L(n)^{-2} n^{-2(2-\ga)} (r/n)^{-2}
\leq C' L(n)^{-2} n^{ 2\ga -3},
\end{align*}
where we used that $\sum_{r\geq n} (r/n)^{-2} \sim n \int_1^{\infty} x^{-2}$ as $n\to\infty$. Therefore, since $\ga<1$, we~get
\[
\sup_{N,M} U_{N,M} \leq \sum_{n=1}^{+\infty} W_n <+\infty\,.
\]
\end{step}

%%\medskip
\begin{step}{Case $\ga > 1$}
First of all, it is immediate that $\sum_{n=1}^N 1/a_n$ diverges as a regularly varying function with exponent $\rho$, since $a_n \sim \psi(n) n^{\frac{1}{\min(\ga,2)}}$, see~\eqref{an}: it directly gives~\eqref{suman}.



We now prove~\eqref{eq:asympUn}. We fix $\gep>0$, and denote, in complement to the definition of $W_n$ above
\[
W_n^{(\gep)} := \sum_{r = 0}^{\lfloor \frac1\gep a_n \rfloor} \bP((n,n+r) \in \tau)^2\,.
\]


As a preliminary, we show that there exists some $n_\gep$ such that, provided that~$n\geq n_{\gep}$
\begin{equation}\label{eq:sumk}
(1-\gep) \frac{c_{\ga}}{a_n} \leq W_n^{(\gep)} \leq W_n \leq (1+\gep) \frac{c_{\ga}}{a_n}.
\end{equation}
Note that we also have that $W_n \leq \sum_{r = 0}^{+\infty} \bP((n,n+r) \in \tau) \leq 1$ for any $n$, since there is at most one renewal in the column $\{n\}\times \bbN$.


To prove~\eqref{eq:sumk}, we use Theorem~\ref{th:reD} to get that uniformly for $0\leq r\leq \tfrac 1\gep a_n$, we have $\bP((n,n+r) \in \tau)^2 \sim (a_n)^{-2} c_{\ga}(r/a_n)^2$ as $n\to\infty$. Hence, provided that $n$ is large enough, we get that
\[
a_n W_n^{(\gep)} \geq (1-\gep^2) \frac{1}{a_n} \sum_{r=0}^{\lfloor\frac1\gep a_n \rfloor} c_{\ga}(r/a_n)^2 \geq (1-2\gep^2) \int_{0}^{1/\gep} c_{\ga}(t)^2 \dd t\,,
\]
the last inequality holding by Riemann--sum approximation. Note that a similar upper bound, with $1-2\gep^2$ replaced with $1+2\gep^2$ holds. Now, thanks to~\eqref{eq:offdiag} (and since $1+\ga-\gd \geq 3/2$), there exists a constant $c>0$ such that
\[
a_n(W_n-W_n^{(\gep)}) = a_n \sum_{r > \frac1\gep a_n}^{+\infty} \bP((n,n+k) \in \tau)^2 \leq c\,\frac{1}{a_n} \sum_{r>\frac1\gep a_n} \Big(\frac{r}{a_n} \Big)^{-3} \leq c' \gep^2,
\]
where the last inequality also comes from a Riemann--sum approximation. Finally, note that $c_\ga -\int_{0}^{1/\gep} c_{\ga}(t)^2 \dd t $ is positive, and thanks to~\eqref{eq:offdiag}, it is smaller than $\int_{1/\gep}^{\infty} c t^{-3} \dd t \leq c'' \gep^2$. In the end, we get that, provided that $n$ is large enough,
\begin{equation}
(1-2\gep^2) (c_{\ga} - c'' \gep^2) \leq a_n W_n^{(\gep)} \leq a_n W_n \leq (1+2\gep^2) c_{\ga} + c' \gep^2,
\end{equation}
which gives~\eqref{eq:sumk} provided that $\gep$ has been fixed small enough.
\end{step}

%%position de la fin de last case )à confirmer.
%%\smallskip
We are now ready to estimate $U_{N,N}$. We write
\[
U_{N,N} = 2\sum_{n=0}^{N} \sum_{r=0}^{N-n} \bP((n,n+r) \in\tau)^2 - \sum_{n=0}^N \bP((n,n)\in\tau)^2.
\]
The second sum is negligible compared to $\sum_{n=1}^N \tfrac{1}{a_n}\sim \varphi(N) N^{\rho}$, since $\bP({(n,n)\in\tau})^2 \sim c (a_n)^{-2}$, with $a_n \to +\infty$. We therefore focus on the first sum.

An upper bound is simply
\[
\sum_{n=0}^N \sum_{r=0}^{N-n} \bP((n,n+r) \in\tau)^2 \leq \sum_{n=0}^N W_n,
\]
and since we have that $W_n \sim c_{\ga}/a_n$ together with~\eqref{suman}, we get that for $n$ large enough
\[
\sum_{n=0}^N \sum_{r=0}^{N-n} \bP((n,n+r) \in\tau)^2 \leq (1+2\gep) c_{\ga} \sum_{n=1}^N \frac{1}{a_n} \leq (1+3\gep) \varphi(N) N^{\rho}\,.
\]
For a lower bound, because $a_N\leq \gep N$ provided that $N$ is large enough, we have
\[
\sum_{n=0}^{N} \sum_{k=0}^{N-n} \bP((n,n+k) \in\tau)^2 \geq \sum_{n=n_{\gep}}^{(1-\gep) N} W_n^{(\gep)} \geq (1-2\gep) c_{\ga} \sum_{n=1}^{(1-\gep)N} \frac{1}{a_n} \geq (1- c\gep)c_{\ga} \varphi(N) N^{\rho},
\]
where we used the lower bound~\eqref{eq:sumk} valid for $n$ large enough, together with~\eqref{suman} for the last inequality.

%%\smallskip
We now turn to estimating $\hat U (\gl)$ as $\gl \downarrow 0$. By symmetry, we can write that
\[
\hat U(\gl) = 2\sum_{n=0}^{+\infty} e^{-2 \gl n}\sum_{r=0}^{+\infty} e^{-\gl r} \bP((n,n+r)\in\tau)^2 - \sum_{n=0}^{+\infty} e^{-2\gl n} \bP((n,n)\in\tau)^2.
\]
The second term is negligible compared to $ \varphi(1\mk/\mk\gl) \gl^{-\rho}$ as $\gl \!\downarrow\! 0$, since $\sum_{n=0}^N\Mk \bP(\mk{(n,\Mk n)\!\in\!\tau})^2$ is negligible compared to $ \varphi(N)N^{\rho}$, by standard properties of Laplace transforms, and we again focus on the first term.

First of all, an upper bound is
\[
\sum_{n=0}^{+\infty} e^{-2 \gl n}\sum_{r=0}^{+\infty} e^{-\gl r} \bP((n,n+r)\in\tau)^2 \leq \sum_{n=0}^{+\infty} e^{-2\gl n} W_n\,.
\]
Since $\sum_{n=0}^{N} W_n \sim c_{\ga} \varphi(N)N^{\rho}$, we get by standard properties of Laplace transforms (see~\cite[Corollary~1.7.3]{cf:BGT}) that
\[
\sum_{n=0}^{+\infty} e^{-2\gl n} W_n \sim \frac{c_\ga}{\Gamma(1+\rho)} \varphi(1/2\gl) (2\gl)^{-\rho} \qquad \text{as } \gl \downarrow 0\,.
\]
For a lower bound, we get that
\[
\sum_{n=0}^{+\infty} e^{-2 \gl n}\sum_{r=0}^{+\infty} e^{-\gl r} \bP((n,n+r)\in\tau)^2 \geq \sum_{n=0}^{+\infty} e^{-2\gl n} e^{-\gl a_n /\gep} W_n^{(\gep)}\,.
\]

Now, we use that there is some $n_{\gep}$ such that for $n\geq n_{\gep}$ we have that $W_n^{(\gep)}\geq (1-\gep) c_{\ga} /a_n$ (see~\eqref{eq:sumk}), and that $a_n/\gep \leq \gep n$. We therefore get that
\begin{align*}
\sum_{n=0}^{+\infty} e^{-2 \gl n}&\sum_{r=0}^{+\infty} e^{-\gl r} \bP((n,n+r)\in\tau)^2 \\
& \geq (1-\gep) c_\ga \sum_{n=n_{\gep}}^{+\infty} e^{-2(1+\gep)\gl n} \frac{1}{a_n} \stackrel{\gl\to 0}{\sim} \frac{(1-\gep) c_\ga}{\Gamma(1+\rho)} \varphi(1/\gl) (2(1+\gep) \gl)^{-\rho}\,,
\end{align*}
where we used again~\cite[Corollary~1.7.3]{cf:BGT} for the last asymptotics.

By letting $\gep \downarrow 0$, we obtain matching upper and lower bound, so that~\eqref{eq:asympUn} is proved.
\end{proofc}

We now use Proposition~\ref{prop:transience}, and in particular the estimate of the Laplace transform $\hat U(\gl)$, to obtain estimates on the tail probability of the intersection renewal $\sigma=\tau\cap \tau'$. More precisely, we define $\underline{\sigma}:= \sigma^{(1)} + \sigma^{(2)}$ and estimate $\bP^{\otimes 2}(\underline \sigma _1 >N)$.

\begin{lemm}\label{lem:tau1}
Assume that $\ga>1$. Then recalling that $\rho=1-\min(\ga,2)^{-1} \in[0,1/2]$, we get that
\begin{equation}
\bP^{\otimes 2} \left(\underline{\sigma}_1 > N \right) \stackrel{N \to \infty}\sim \frac{2^{\rho} \sin(\pi \rho)}{ \pi \rho} \ (U_{N,N})^{-1} \stackrel{N \to \infty}\sim C_{\ga,\rho}\,\varphi(N)^{-1} N^{-\rho}\,.
\end{equation}
\end{lemm}

\begin{proof}
Recall the definition of $\hat U (\gl) = \sum_{n,m \ge 0} e^{- \lambda (n+m)} \bP^{\otimes 2} \left((n,m) \in \sigma \right) $. We also set, for any $\gl>0$,
\begin{equation}
\hat{K}(\lambda) := \sum_{n,m \ge 1} e^{- \lambda (n+m)} \bP^{\otimes 2} (\sigma_1 = (n,m)) = \sum_{k \ge 2} e^{- \lambda k} \bP^{\otimes 2} (\underline{\sigma}_1 = k)\,.
\end{equation}
The key idea of this proof is the following identity
\begin{equation}\label{eq:id}
\hat U(\gl) = 1+ \hat K(\gl) \hat U(\gl) \quad\Longleftrightarrow \quad 1- \hat{K}(\lambda) = \frac{1}{\hat{U}(\lambda)}\,,
\end{equation}
which is obtained from the identity
\begin{equation}
\bP^{\otimes 2} \left((n,m) \in \sigma \right) = \unbf_{\{n=m=0\}} + \sum_{i=1}^{n}\sum_{j=1}^{m} \bP^{\otimes 2} (\sigma_1 = (i,j)) \bP^{\otimes 2} ((n-i,n-j) \in\sigma).\!
\end{equation}

Now, since we know the behavior of $\hat U(\gl)$ as $\gl\downarrow0$, we get the behavior of $\hat K(\gl)$, from which we should be able to infer that of $\bP^{\otimes 2}(\underline \sigma_1 >N)$. Let us develop here how we proceed: we use~\cite[Corollary~1.7.3 and Theorem~8.7.3]{cf:BGT}. We can view $\underline\sigma$ as a renewal process with inter-arrival distribution $\bP^{\otimes 2} (\underline \sigma_1=k) = \bP^{\otimes 2} (\sigma_1^{(1)} + \sigma_1^{(2)}=k)$, and we set $u_n:= \bP^{\otimes 2} (n \in \underline \sigma)$ its renewal mass function, so we have $\hat U(\gl) = \sum_{n=0}^{\infty} e^{-\gl n} u_n$ (and~\eqref{eq:id} is standard from the one-dimensional renewal equation). Now, \cite[Corollary~1.7.3]{cf:BGT} tells that since $\hat U(\gl)$ is regularly varying with exponent $-\rho$ (recall $\rho = 1-\min(\ga,2)^{-1}$), we have that $\sum_{n=0}^N u_n \sim \Gamma(1+\rho) \hat U(1/N) \sim 2^{-\rho} U_{N,N}$ (where we used~\eqref{eq:asympUn}). In turn~\cite[Theorem~8.7.3]{cf:BGT} gives that
\[
\bP^{\otimes 2} (\underline\sigma _1 >N) \stackrel{N\to\infty}{\sim} \frac{(2^{-\rho} U_{N,N})^{-1}}{\Gamma(1+\rho) \Gamma(1-\rho)},
\]
and we are done.
\end{proof}

\begin{rema}\label{rem:alpha=1}
The case $\ga=1$ has been left aside, mostly to avoid technicalities. Denote $\mu(n):=\bE[\min(\tau_1^{(1)}, n)]$ the truncated first moment of $\tau_1^{(1)}$. It is shown in~\cite[Theorem~3.4]{cf:B} (or~\eqref{eq:zfb} in the symmetric context) that along the favorite direction, for $n\to \infty$ and $r$ with $r/a_{n/\mu(n)} \to t\in \bbR_+$, we have
\begin{equation}\label{alpha=1}
\bP((n,n+r)\in\tau) \sim \frac{c_1(t)}{\mu(n) a_{n/\mu(n)}}\,,
\end{equation}
with $c_1(t):= \int_{-\infty}^{+\infty} g_{\alpha}(x,(1+t)x)\, \dd x$. Notice that $n/\mu(n)$ is the typical number of steps to reach distance $n$. Again, estimates away from the favorite direction are provided in~\cite[Theorem~4.2]{cf:B} (or~\eqref{eq:rho} in the symmetric case): for any $\gd>0$, there is a constant $C_{\gd}$ such that for any $r \ge a_{n/\mu(n)}$,
\begin{equation}\label{alpha=1away}
\bP((n,n+r)\in\tau) \le \frac{C_{\gd}}{\mu(n) a_{n/\mu(n)}}\,\Big(\frac{r}{\mu(n) a_{n/\mu(n)}} \Big)^{-2+\gd}\,.
\end{equation}
This shows that the main contribution to $U_{N,N}$ comes also here from the terms close to the diagonal, that is
\begin{align*}
U_{N,N} &\asymp 2\sum_{n=1}^N \sum_{r=0}^{a_{n/\mu(n)}} \bP((n,n+r)\in\tau)^2 \asymp \sum_{n=1}^N \frac{1}{\mu(n)^2 a_{n/\mu(n)}}\,.
\end{align*}

(We denoted $x_n\asymp y_n$ if $x_n/y_n$ is bounded away from $0$ and $+\infty$.) Let us stress that we have $\mu(n) \sim \mu(a_{n/\mu(n)})$ (this comes from~\cite[Lemma~4.3]{cf:B17}): by a change of variable $x= n/\mu(n)$ (comparing the sum to an integral, and considering $\mu(n),a_n$ as functions of positive real numbers), we get that
\[
U_{N,N}\asymp \int_{1}^{N/\mu(N)} \frac{\dd x}{ a_{x} \mu(a_x)} \asymp \int_1^{a_{N/\mu(N)}} \frac{\dd u}{u L(u) \mu(u)}\,,
\]
where we used another change of variables $u \!=\!a_x$ ($\dd x \!\sim\! L(u)^{-1}\dd u$, since $n\!\sim\! a_n/L(a_n)$). As a conclusion, we expect to have the following criterion:
\[
\sigma=\tau\cap \tau' \text{ is persistent } \;\; \Longleftrightarrow \;\; \sum_{n\geq 1} \frac{1}{a_n \mu(a_n)} =+\infty \;\; \Longleftrightarrow \;\; \sum_{n\geq 1} \frac{1}{n L(n)\mu(n)} =+\infty.
\]
As an example, if $L(n) =(\log n)^{\kappa}$ with $\kappa \geq -1$, then $\mu(n) \sim c_{\kappa} \max(1, (\log n)^{1+\kappa})$ and hence $\sigma=\tau\cap \tau'$ should be persistent if and only if $\kappa>0$.
\end{rema}
\end{appendix}

